\documentclass[../../main.tex]{subfiles}
\begin{document}

{\bf [解]}

与えられた二次方程式$x^2-4x-1=0$の解と係数の関係から
\begin{align}
  \alpha+\beta=4, \alpha\beta=-1 \label{eq:1}
\end{align}
である．

\paragraph{(1)}

$S_1,S_2,S_3$を順番に計算すると
\begin{align}
  S_1 = \alpha+\beta = 4
\end{align}
\begin{align}
  S_2 = \alpha^2+\beta^2 = (\alpha+\beta)^2-2\alpha\beta = 18
\end{align}
\begin{align}
  S_3 = (\alpha+\beta)^3-3\alpha\beta(\alpha+\beta) = 76
\end{align}
だから
\begin{align}
  S_1 = 4, S_2 = 18, S_3 = 76
\end{align}
である．

一方で$n \ge 3$の時，
\begin{align}
  S_n
  &= \alpha^n+\beta^n \\
  &= (\alpha+\beta)\left(\alpha^{n-1}+\beta^{n-1}\right) - \alpha\beta\left(\alpha^{n-2}+\beta^{n-2}\right) \\
  &= 4S_{n-1} + S_{n-2}
\end{align}
である．

\paragraph{(2)}

$\beta$は与えられた二次方程式の小さい方の解だから
\begin{align}
  \beta^2-4\beta-1=0
\end{align}
である．よって
\begin{align}
  \beta = 2-\sqrt{5}
\end{align}
である．
よって$-1 < \beta < 0$ より $-1 < \beta^3 < 0$，つまりもとめるのは $-1$ である．

\paragraph{(3)}

(1)の$S_1,S_2,S_3$および漸化式より，
$S_n \in \mathbb{N}$ でかつ下1桁は $\{4, 8, 6, 2\}$ の周期数列で $S_{2003}$ の下1ケタは $6$ である．
又(2)と同様に $-1 < \beta^{2003} < 0$ だから，$\alpha^{2003} = S_{2003} - \beta^{2003}$ 以下の最大の整数の下1ケタは $6$ である．

\end{document}
